% Generated draft: requires mathematical and source review.
\documentclass[11pt]{article}
\usepackage[margin=25mm]{geometry}
\usepackage{fontspec}
\setmainfont{Latin Modern Roman}
\usepackage{amsmath,amssymb,mathtools,unicode-math}
\setmathfont{Latin Modern Math}
\usepackage{enumitem,needspace,xurl,hyperref}
\setlength{\parindent}{0pt}
\setlength{\parskip}{5pt}
\setlength{\emergencystretch}{3em}
\newcommand{\N}{\mathbb{N}}
\newcommand{\Z}{\mathbb{Z}}
\newcommand{\Q}{\mathbb{Q}}
\newcommand{\R}{\mathbb{R}}
\newcommand{\C}{\mathbb{C}}
\newcommand{\F}{\mathbb{F}}
\newcommand{\A}{\mathbb{A}}
\newcommand{\PP}{\mathbb{P}}
\newcommand{\E}{\mathbb{E}}
\newcommand{\eps}{\varepsilon}
\newcommand{\dd}{\,\mathrm{d}}
\begin{document}
\section*{UnsolvedMath ID: 1101205}
\textbf{Source problem code:} catalogue entry
\subsection{Definitions and mathematical statement}
Use the surface domain in the supplied research summary: \(S=S_{g,p}\) is a connected orientable surface of finite type, genus \(g\), with \(p\) punctures and no boundary, and \(\xi(S)=3g-3+p\ge1\). Let \(\mathcal T(S)\) be its Teichm\"uller space, with Weil--Petersson distance \(d_{\mathrm{WP}}\), and let \(\operatorname{Mod}(S)\) be its mapping class group with a word metric from a finite generating set. For a metric space \(X\), define its quasi-flat rank as the supremum of integers \(k\ge0\) for which there is a map \(q:\R^k\to X\) and constants \(L\ge1,A\ge0\) satisfying
\[
L^{-1}\|u-v\|-A\le d(q(u),q(v))\le L\|u-v\|+A\qquad(u,v\in\R^k).
\]
Brock asks for these ranks for \((\mathcal T(S),d_{\mathrm{WP}})\) and \(\operatorname{Mod}(S)\). Under the supplied domain and meaning of rank, the research summary records the answers
\[
\operatorname{rank}\operatorname{Mod}(S)=3g-3+p,
\qquad
\operatorname{rank}(\mathcal T(S),d_{\mathrm{WP}})
=\left\lfloor\frac{3g+p-2}{2}\right\rfloor
=\left\lfloor\frac{\xi(S)+1}{2}\right\rfloor.
\]
It identifies the first value with maximal free-abelian subgroup rank and the second with the analogous pants-graph rank, using Brock's quasi-isometry theorem. It describes the second value in terms of disjoint nonannular essential subsurfaces; the relevant subsurfaces must have positive complexity for that interpretation, rather than including pairs of pants. The literal source question does not specify the surface or rank convention; the formulas here retain precisely the \(\xi(S)\ge1\) domain of the summary, without extending them to excluded low-complexity surfaces or surfaces with boundary.

The summary also asserts that maximal quasi-flats are Hausdorff-close to unions of ``Dehn twist flats''/standard orthants. Its intended class of standard pieces and hypotheses are not specified, so that ancillary claim is not promoted to a precise additional theorem here. The summary ends in a truncated consistency-check sentence.
\subsection{Short English statement}
Determine the largest dimension of a Euclidean space that quasi-isometrically embeds into the mapping class group or into Teichm\"uller space with the Weil--Petersson metric. The supplied summary gives the displayed ranks for connected punctured orientable surfaces of complexity at least one.
\subsection{Sources}
\begin{itemize}
\item Local UnsolvedMath registry, permanent ID \(1101205\), exact source code \texttt{catalogue entry}, ``Questions in Geometric Group Theory --- Q 12.5.''
\item Bestvina, \emph{Questions in Geometric Group Theory} (2004), Question 12.5, PDF page 20; attributed to Brock. \url{https://www.math.utah.edu/~bestvina/eprints/questions-updated.pdf}; manual transcription; recorded access 2026-07-29.
\item Supplied research summary attributes maximal abelian rank to Birman--Lubotzky--McCarthy; the Brock--Farb Rank Conjecture to Behrstock--Minsky (2008) and independently Hamenst\"adt; further proofs to Eskin--Masur--Rafi (2017) and Bowditch. It attributes the quasi-flat description to Behrstock--Hagen--Sisto.
\item The same summary cites Brock's theorem, Eskin--Masur--Rafi, Corollary C, earlier Behrstock--Minsky computations, and Bowditch (2020) for Weil--Petersson rank. Complete bibliographic references are not supplied.
\end{itemize}
\end{document}
